\section{电子--声子自能与 Migdal 定理}

有了 Matsubara 图规则，可以把 Fröhlich 哈密顿量 \eqref{eq:frohlich} 推进到可计算的自能（Mahan 第 7 章）。

\subsection{最低阶电子自能}

\begin{equation}
  \Sigma(\mathbf{k},\mathrm{i}\omega_n)
  =\frac{1}{\beta L^d}
  \sum_{\mathbf{q},m\lambda}
  |M_{\mathbf{q}\lambda}|^2
  \mathcal{D}_0(\mathbf{q},\mathrm{i}\nu_m)
  \mathcal{G}(\mathbf{k}-\mathbf{q},\mathrm{i}\omega_n-\mathrm{i}\nu_m),
\label{eq:Sigma_ep}
\end{equation}
$\mathcal{D}_0$ 见 \eqref{eq:D0}。对 $m$ 做玻色频率求和后，实频
\begin{equation}
  \Sigma^R(\mathbf{k},\omega)
  =\sum_{\mathbf{q}\lambda}
  |M_{\mathbf{q}\lambda}|^2
  \left[
  \frac{n_B(\omega_{\mathbf{q}})+1-n_F(\xi_{\mathbf{k}-\mathbf{q}})}
  {\omega-\xi_{\mathbf{k}-\mathbf{q}}-\omega_{\mathbf{q}}+\mathrm{i}0^+}
  +\frac{n_B(\omega_{\mathbf{q}})+n_F(\xi_{\mathbf{k}-\mathbf{q}})}
  {\omega-\xi_{\mathbf{k}-\mathbf{q}}+\omega_{\mathbf{q}}+\mathrm{i}0^+}
  \right].
\label{eq:Sigma_ep_R}
\end{equation}
两项分别对应发射与吸收声子。

\subsection{质量重整与 \texorpdfstring{$\lambda$}{lambda}}

费米面、$\omega\to0$：
\begin{equation}
  \mathrm{Re}\,\Sigma(\omega)\approx-\lambda\omega,
  \qquad
  \lambda=2\int_0^{\infty}\frac{\mathrm{d}\omega}{\omega}\,\alpha^2F(\omega),
\label{eq:lambda}
\end{equation}
Eliashberg 函数 $\alpha^2F$ 为费米面平均的 $|M|^2$ 乘声子谱。于是
\begin{equation}
  \frac{m^*}{m}=1+\lambda.
\end{equation}
$\mathrm{Im}\,\Sigma$ 在 Debye 窗口内给出有限寿命，窗口外迅速变小（光学实验中的“声子结构”）。

\subsection{Migdal 定理}

顶点修正 $\sim\omega_D/E_F\ll1$。金属中可忽略交叉声子线，只保留 \eqref{eq:Sigma_ep} 的彩虹图。$\lambda$ 很大或 $\omega_D\sim E_F$ 时定理失效。

\subsection{弱耦合 Fröhlich 极化子}

单一电子、无费米海时，\eqref{eq:Sigma_ep_R} 在 $T=0$、$k=0$ 给出
\begin{equation}
  \Delta E=-\sum_{\mathbf{q}}\frac{|M_{\mathbf{q}}|^2}{\omega_{\mathbf{q}}+\varepsilon_{\mathbf{q}}},
\end{equation}
极性光学模积出 $\Delta E=-\alpha\omega_{\mathrm{LO}}$，$m^*/m=1+\alpha/6+\cdots$。强耦合需独立玻色子/LLP 变分，见上一节。
